\documentclass[a4paper,12pt]{article}
\usepackage[utf8]{inputenc}
\usepackage[ngerman]{babel}
\usepackage{amsmath}
\usepackage{parskip}

\setlength{\parindent}{0pt}

\title{Modul 293 - Aufgabenblatt 13}
\author{von Lukas Meier}
\date{Unterricht vom 31.03.2026}

\begin{document}

\maketitle

\section*{Aufgabenblatt: Responsive Grundlagen}

\subsection*{Teil 1: Begriffe verstehen}
Beantworte die folgenden Fragen kurz schriftlich:
\begin{itemize}
  \item Was bedeutet \emph{responsive} bei einer Webseite?
  \item Wozu dient die Meta-Zeile \texttt{viewport} im \texttt{head}?
  \item Warum sind fixe Pixelbreiten auf kleinen Geräten oft problematisch?
\end{itemize}

\subsection*{Teil 2: Media Queries lesen}
Betrachte die folgende CSS-Regel:

\begin{verbatim}
@media (max-width: 700px) {
  nav a {
    display: block;
  }
}
\end{verbatim}

Bearbeite dazu die Fragen:
\begin{itemize}
  \item Wann gilt diese Regel?
  \item Was verändert sich bei den Links im \texttt{nav}?
  \item Für welche Geräte könnte diese Regel sinnvoll sein?
\end{itemize}

\subsection*{Teil 3: max-width und min-width}
Erkläre den Unterschied zwischen:
\begin{itemize}
  \item \texttt{@media (max-width: 700px)}
  \item \texttt{@media (min-width: 700px)}
\end{itemize}

Schreibe zu beiden Varianten je einen kurzen Beispielsatz in eigenen Worten.

\subsection*{Teil 4: Prozentbreiten anwenden}
Erstelle eine einfache HTML-Seite mit einem Container \texttt{wrapper} und zwei Bereichen:
\begin{itemize}
  \item \texttt{left}
  \item \texttt{right}
\end{itemize}

Vorgaben:
\begin{itemize}
  \item Der Container soll z.\,B. \texttt{width: 90\%} haben.
  \item Beide Bereiche sollen auf einem breiten Bildschirm nebeneinander stehen.
  \item Beide Bereiche sollen ungefähr gleich breit sein.
\end{itemize}

\subsection*{Teil 5: Layout für kleine Bildschirme anpassen}
Erweitere deine Lösung aus Teil 4 mit einer Media Query:
\begin{itemize}
  \item Bei kleiner Bildschirmbreite sollen \texttt{left} und \texttt{right} untereinander dargestellt werden.
  \item Beide Bereiche sollen dann \texttt{width: 100\%} erhalten.
\end{itemize}

Teste dein Ergebnis, indem du das Browserfenster schmaler und breiter ziehst.

\subsection*{Teil 6: Code ergänzen}
Vervollständige den folgenden CSS-Code:

\begin{verbatim}
.card {
  width: ______;
  display: inline-block;
}

@media (max-width: 700px) {
  .card {
    width: ______;
    display: ______;
  }
}
\end{verbatim}

Ziel:
\begin{itemize}
  \item Auf grossen Bildschirmen soll die Karte nur einen Teil der Breite belegen.
  \item Auf kleinen Bildschirmen soll sie die ganze Breite nutzen.
\end{itemize}

\subsection*{Teil 7: Best Practices reflektieren}
Beurteile die folgenden Aussagen mit \glqq sinnvoll\grqq{} oder \glqq eher ungünstig\grqq{} und begründe kurz:
\begin{itemize}
  \item Ein Hauptbereich bekommt immer \texttt{width: 1200px}.
  \item Buttons und Texte werden auf dem Handy zusätzlich getestet.
  \item Ein Container bekommt \texttt{width: 90\%} und \texttt{margin: 0 auto}.
  \item Das Layout wird nur auf einem grossen Monitor angeschaut.
\end{itemize}

\subsection*{Teil 8: Praxisaufgabe}
Erstelle eine kleine Beispielseite zum Thema \glqq Mein Wochenplan\grqq{} oder \glqq Mein Lieblingsort\grqq{}.

Anforderungen:
\begin{itemize}
  \item eine HTML-Datei und eine CSS-Datei
  \item ein Container mit prozentualer Breite
  \item mindestens zwei Inhaltsbereiche
  \item mindestens eine Media Query mit \texttt{max-width} oder \texttt{min-width}
  \item sichtbarer Unterschied zwischen schmalem und breitem Layout
\end{itemize}

\subsection*{Teil 9: Reflexion}
Beantworte die folgenden Fragen kurz:
\begin{itemize}
  \item Was war heute bei Responsive Design neu für dich?
  \item Was hat mit Media Queries bereits gut funktioniert?
  \item Wo brauchst du noch Übung?
\end{itemize}

\end{document}
